% --- Archon physics formalization source begin ---
% archon:physics
% archon:covers PhyXMiniProblems/problem_phyx_mini_0900.lean
% archon:source-report reports/phyx_mini/problem_phyx_mini_0900.source.json

\chapter{Physics problem phyx\_mini\_0900}
\label{ch:PhyXMiniProblems_problem_phyx_mini_0900}

\paragraph{Problem source.}
\#\# Physical scenario

Charge \textbackslash{}q\_2 is in static equilibrium.

\#\# Question

What is \textbackslash{}q\_1?

\#\# Answer choices

- A: A: -12nC

- B: B: 20nC

- C: C: -20nC

- D: D: 5nC

\#\# Figure caption (auxiliary; use the image as primary evidence)

The image depicts a linear arrangement of three circular objects. Here's a detailed description:

- **Left Circle**: Labeled as \textbackslash{}( q\_1 \textbackslash{}), this circle is light blue/green in color.

- **Middle Circle**: Labeled as "5.0 nC" with a positive sign (“+”) inside it, this circle is red in color. It is centered between two arrows pointing to the left and right.

- **Right Circle**: Labeled as \textbackslash{}( q\_2 \textbackslash{}), this circle is a light red color.

- **Distances**: The distance between the left circle (\textbackslash{}( q\_1 \textbackslash{})) and the middle circle (5.0 nC) is marked as "10 cm." Similarly, the distance between the middle circle and the right circle (\textbackslash{}( q\_2 \textbackslash{})) is also "10 cm."

- **Arrows and Lines**: There is a horizontal line with double-headed arrows extending between the objects, indicating the distances between them.

\#\# Dataset metadata

Electrostatics; Physical Model Grounding Reasoning, Multi-Formula Reasoning

\paragraph{Recorded answer/context.}
C: C: -20nC

\paragraph{Figure/image path.}
/root/proposal\_for\_physic/science-mango/phyx\_mini\_run/phyx\_data/test\_image/900.png

\paragraph{Formalization target.}
create a compiling Lean file with sorry bodies at `PhyXMiniProblems/problem\_phyx\_mini\_0900.lean`. The Lean declarations must preserve the physical quantities, dimensions or dimensional roles, figure labels, governing-law hypotheses, and final relation expressed by this problem.
Use Mathlib/Physlib names found through LeanExplore where available. If a physics API is missing, introduce faithful local abstractions rather than scalar placeholder aliases.

\begin{theorem}[Physics formalization target]
\label{thm:physics:phyx_mini_0900:target}
\lean{PhyXMiniProblems.ProblemPhyXMini0900.problem_phyx_mini_0900}
\uses{def:physics:phyx-mini-0900:phyxminiproblems-problemphyxmini0900-chargeinnanocoulombs, def:physics:phyx-mini-0900:phyxminiproblems-problemphyxmini0900-collinearthreechargeequilibriumsetup, def:physics:phyx-mini-0900:phyxminiproblems-problemphyxmini0900-matchesthreepointchargeequilibriumscenario, def:physics:phyx-mini-0900:phyxminiproblems-problemphyxmini0900-matchessuppliedcollinearchargefigure, def:physics:phyx-mini-0900:phyxminiproblems-problemphyxmini0900-hasphysicalthreechargeparameters, def:physics:phyx-mini-0900:phyxminiproblems-problemphyxmini0900-satisfiesaxialpointchargecoulombforcelaw, def:physics:phyx-mini-0900:phyxminiproblems-problemphyxmini0900-satisfiesforcesuperpositiononq2, def:physics:phyx-mini-0900:phyxminiproblems-problemphyxmini0900-satisfiesstaticequilibriumforcebalance}
The assigned autoformalize agent should translate this physics problem into Lean declarations in the covered file, with theorem and lemma proof bodies written as `by sorry`.
\end{theorem}
\begin{proof}
This is an autoformalization task, not a proof task. Produce statements that can later be proved by the physics prover without weakening the physical model.
\end{proof}

% --- Archon declaration topology begin ---
\section{Declaration topology}
\begin{definition}[force Dimension]
  \label{def:physics:phyx-mini-0900:phyxminiproblems-problemphyxmini0900-forcedimension}
  \lean{PhyXMiniProblems.ProblemPhyXMini0900.forceDimension}
  The physical dimension M L T⁻² of a force component
\end{definition}
\begin{proof}
  This is the declared model definition of force Dimension.
\end{proof}
\begin{definition}[Signed Charge Quantity]
  \label{def:physics:phyx-mini-0900:phyxminiproblems-problemphyxmini0900-signedchargequantity}
  \lean{PhyXMiniProblems.ProblemPhyXMini0900.SignedChargeQuantity}
  A signed, unit-independent physical electric charge
\end{definition}
\begin{proof}
  This is the declared model definition of Signed Charge Quantity.
\end{proof}
\begin{definition}[Axial Position Quantity]
  \label{def:physics:phyx-mini-0900:phyxminiproblems-problemphyxmini0900-axialpositionquantity}
  \lean{PhyXMiniProblems.ProblemPhyXMini0900.AxialPositionQuantity}
  A signed, unit-independent position coordinate on the horizontal axis
\end{definition}
\begin{proof}
  This is the declared model definition of Axial Position Quantity.
\end{proof}
\begin{definition}[Length Quantity]
  \label{def:physics:phyx-mini-0900:phyxminiproblems-problemphyxmini0900-lengthquantity}
  \lean{PhyXMiniProblems.ProblemPhyXMini0900.LengthQuantity}
  A nonnegative, unit-independent physical separation
\end{definition}
\begin{proof}
  This is the declared model definition of Length Quantity.
\end{proof}
\begin{definition}[Axial Force Component Quantity]
  \label{def:physics:phyx-mini-0900:phyxminiproblems-problemphyxmini0900-axialforcecomponentquantity}
  \lean{PhyXMiniProblems.ProblemPhyXMini0900.AxialForceComponentQuantity}
  \uses{def:physics:phyx-mini-0900:phyxminiproblems-problemphyxmini0900-forcedimension}
  A signed, unit-independent horizontal force component
\end{definition}
\begin{proof}
  This is the declared model definition of Axial Force Component Quantity.
\end{proof}
\begin{definition}[charge Readout]
  \label{def:physics:phyx-mini-0900:phyxminiproblems-problemphyxmini0900-chargereadout}
  \lean{PhyXMiniProblems.ProblemPhyXMini0900.chargeReadout}
  \uses{def:physics:phyx-mini-0900:phyxminiproblems-problemphyxmini0900-signedchargequantity}
  Read a physical charge in a selected Physlib charge unit
\end{definition}
\begin{proof}
  This is the declared model definition of charge Readout.
\end{proof}
\begin{definition}[nanocoulomb Unit]
  \label{def:physics:phyx-mini-0900:phyxminiproblems-problemphyxmini0900-nanocoulombunit}
  \lean{PhyXMiniProblems.ProblemPhyXMini0900.nanocoulombUnit}
  The charge unit 10⁻⁹ C used by the middle label and answer choices
\end{definition}
\begin{proof}
  This is the declared model definition of nanocoulomb Unit.
\end{proof}
\begin{definition}[charge In Coulombs]
  \label{def:physics:phyx-mini-0900:phyxminiproblems-problemphyxmini0900-chargeincoulombs}
  \lean{PhyXMiniProblems.ProblemPhyXMini0900.chargeInCoulombs}
  \uses{def:physics:phyx-mini-0900:phyxminiproblems-problemphyxmini0900-signedchargequantity, def:physics:phyx-mini-0900:phyxminiproblems-problemphyxmini0900-chargereadout}
  Coherent-SI coulomb readout of a signed charge
\end{definition}
\begin{proof}
  This is the declared model definition of charge In Coulombs.
\end{proof}
\begin{definition}[charge In Nanocoulombs]
  \label{def:physics:phyx-mini-0900:phyxminiproblems-problemphyxmini0900-chargeinnanocoulombs}
  \lean{PhyXMiniProblems.ProblemPhyXMini0900.chargeInNanocoulombs}
  \uses{def:physics:phyx-mini-0900:phyxminiproblems-problemphyxmini0900-signedchargequantity, def:physics:phyx-mini-0900:phyxminiproblems-problemphyxmini0900-chargereadout, def:physics:phyx-mini-0900:phyxminiproblems-problemphyxmini0900-nanocoulombunit}
  Nanocoulomb readout used by the printed charge label
\end{definition}
\begin{proof}
  This is the declared model definition of charge In Nanocoulombs.
\end{proof}
\begin{definition}[length Readout]
  \label{def:physics:phyx-mini-0900:phyxminiproblems-problemphyxmini0900-lengthreadout}
  \lean{PhyXMiniProblems.ProblemPhyXMini0900.lengthReadout}
  \uses{def:physics:phyx-mini-0900:phyxminiproblems-problemphyxmini0900-lengthquantity}
  Read a nonnegative physical separation in a selected length unit
\end{definition}
\begin{proof}
  This is the declared model definition of length Readout.
\end{proof}
\begin{definition}[length In Meters]
  \label{def:physics:phyx-mini-0900:phyxminiproblems-problemphyxmini0900-lengthinmeters}
  \lean{PhyXMiniProblems.ProblemPhyXMini0900.lengthInMeters}
  \uses{def:physics:phyx-mini-0900:phyxminiproblems-problemphyxmini0900-lengthquantity, def:physics:phyx-mini-0900:phyxminiproblems-problemphyxmini0900-lengthreadout}
  Coherent-SI metre readout of a separation
\end{definition}
\begin{proof}
  This is the declared model definition of length In Meters.
\end{proof}
\begin{definition}[length In Centimeters]
  \label{def:physics:phyx-mini-0900:phyxminiproblems-problemphyxmini0900-lengthincentimeters}
  \lean{PhyXMiniProblems.ProblemPhyXMini0900.lengthInCentimeters}
  \uses{def:physics:phyx-mini-0900:phyxminiproblems-problemphyxmini0900-lengthquantity, def:physics:phyx-mini-0900:phyxminiproblems-problemphyxmini0900-lengthreadout}
  Centimetre readout used by the two dimension arrows
\end{definition}
\begin{proof}
  This is the declared model definition of length In Centimeters.
\end{proof}
\begin{definition}[position In Meters]
  \label{def:physics:phyx-mini-0900:phyxminiproblems-problemphyxmini0900-positioninmeters}
  \lean{PhyXMiniProblems.ProblemPhyXMini0900.positionInMeters}
  \uses{def:physics:phyx-mini-0900:phyxminiproblems-problemphyxmini0900-axialpositionquantity}
  Coherent-SI metre readout of a signed axial position
\end{definition}
\begin{proof}
  This is the declared model definition of position In Meters.
\end{proof}
\begin{definition}[force Component In Newtons]
  \label{def:physics:phyx-mini-0900:phyxminiproblems-problemphyxmini0900-forcecomponentinnewtons}
  \lean{PhyXMiniProblems.ProblemPhyXMini0900.forceComponentInNewtons}
  \uses{def:physics:phyx-mini-0900:phyxminiproblems-problemphyxmini0900-axialforcecomponentquantity}
  Coherent-SI newton readout of a signed axial force component
\end{definition}
\begin{proof}
  This is the declared model definition of force Component In Newtons.
\end{proof}
\begin{lemma}[charge In Nanocoulombs eq scaled charge In Coulombs]
  \label{lem:physics:phyx-mini-0900:phyxminiproblems-problemphyxmini0900-chargeinnanocoulombs-eq-scaled-chargeincoulombs}
  \lean{PhyXMiniProblems.ProblemPhyXMini0900.chargeInNanocoulombs_eq_scaled_chargeInCoulombs}
  \uses{def:physics:phyx-mini-0900:phyxminiproblems-problemphyxmini0900-signedchargequantity, def:physics:phyx-mini-0900:phyxminiproblems-problemphyxmini0900-chargeincoulombs, def:physics:phyx-mini-0900:phyxminiproblems-problemphyxmini0900-chargeinnanocoulombs}
  The Physlib named units make these scale conversions consequences rather than extra physical hypotheses. They are recorded as helper lemmas for the later algebraic proof
\end{lemma}
\begin{proof}
  The conclusion follows from the governing relations and figure readouts named in the statement of charge In Nanocoulombs eq scaled charge In Coulombs.
\end{proof}
\begin{lemma}[length In Centimeters eq scaled length In Meters]
  \label{lem:physics:phyx-mini-0900:phyxminiproblems-problemphyxmini0900-lengthincentimeters-eq-scaled-lengthinmeters}
  \lean{PhyXMiniProblems.ProblemPhyXMini0900.lengthInCentimeters_eq_scaled_lengthInMeters}
  \uses{def:physics:phyx-mini-0900:phyxminiproblems-problemphyxmini0900-lengthquantity, def:physics:phyx-mini-0900:phyxminiproblems-problemphyxmini0900-lengthinmeters, def:physics:phyx-mini-0900:phyxminiproblems-problemphyxmini0900-lengthincentimeters}
  This lemma records the length In Centimeters eq scaled length In Meters component of the problem-specific physical model
\end{lemma}
\begin{proof}
  The conclusion follows from the governing relations and figure readouts named in the statement of length In Centimeters eq scaled length In Meters.
\end{proof}
\begin{definition}[Charge Site]
  \label{def:physics:phyx-mini-0900:phyxminiproblems-problemphyxmini0900-chargesite}
  \lean{PhyXMiniProblems.ProblemPhyXMini0900.ChargeSite}
  The three circular charges, named by their left-to-right figure roles
\end{definition}
\begin{proof}
  This is the declared model definition of Charge Site.
\end{proof}
\begin{definition}[Adjacent Gap]
  \label{def:physics:phyx-mini-0900:phyxminiproblems-problemphyxmini0900-adjacentgap}
  \lean{PhyXMiniProblems.ProblemPhyXMini0900.AdjacentGap}
  The two labelled adjacent gaps in the image
\end{definition}
\begin{proof}
  This is the declared model definition of Adjacent Gap.
\end{proof}
\begin{definition}[Figure Charge Color]
  \label{def:physics:phyx-mini-0900:phyxminiproblems-problemphyxmini0900-figurechargecolor}
  \lean{PhyXMiniProblems.ProblemPhyXMini0900.FigureChargeColor}
  Fill colours visibly distinguishing the three charge circles
\end{definition}
\begin{proof}
  This is the declared model definition of Figure Charge Color.
\end{proof}
\begin{definition}[Figure Sign Glyph]
  \label{def:physics:phyx-mini-0900:phyxminiproblems-problemphyxmini0900-figuresignglyph}
  \lean{PhyXMiniProblems.ProblemPhyXMini0900.FigureSignGlyph}
  The only sign glyph visible inside a charge circle is a plus
\end{definition}
\begin{proof}
  This is the declared model definition of Figure Sign Glyph.
\end{proof}
\begin{definition}[expected Charge Text]
  \label{def:physics:phyx-mini-0900:phyxminiproblems-problemphyxmini0900-expectedchargetext}
  \lean{PhyXMiniProblems.ProblemPhyXMini0900.expectedChargeText}
  \uses{def:physics:phyx-mini-0900:phyxminiproblems-problemphyxmini0900-chargesite}
  Literal charge text expected above each of the three circles
\end{definition}
\begin{proof}
  This is the declared model definition of expected Charge Text.
\end{proof}
\begin{definition}[expected Charge Color]
  \label{def:physics:phyx-mini-0900:phyxminiproblems-problemphyxmini0900-expectedchargecolor}
  \lean{PhyXMiniProblems.ProblemPhyXMini0900.expectedChargeColor}
  \uses{def:physics:phyx-mini-0900:phyxminiproblems-problemphyxmini0900-chargesite, def:physics:phyx-mini-0900:phyxminiproblems-problemphyxmini0900-figurechargecolor}
  Literal circle colour expected for each charge site
\end{definition}
\begin{proof}
  This is the declared model definition of expected Charge Color.
\end{proof}
\begin{definition}[expected Sign Glyph]
  \label{def:physics:phyx-mini-0900:phyxminiproblems-problemphyxmini0900-expectedsignglyph}
  \lean{PhyXMiniProblems.ProblemPhyXMini0900.expectedSignGlyph}
  \uses{def:physics:phyx-mini-0900:phyxminiproblems-problemphyxmini0900-chargesite, def:physics:phyx-mini-0900:phyxminiproblems-problemphyxmini0900-figuresignglyph}
  Sign glyph expected inside each circle
\end{definition}
\begin{proof}
  This is the declared model definition of expected Sign Glyph.
\end{proof}
\begin{definition}[expected Printed Charge Nanocoulombs]
  \label{def:physics:phyx-mini-0900:phyxminiproblems-problemphyxmini0900-expectedprintedchargenanocoulombs}
  \lean{PhyXMiniProblems.ProblemPhyXMini0900.expectedPrintedChargeNanocoulombs}
  \uses{def:physics:phyx-mini-0900:phyxminiproblems-problemphyxmini0900-chargesite}
  Optional numerical nanocoulomb label attached to each charge site
\end{definition}
\begin{proof}
  This is the declared model definition of expected Printed Charge Nanocoulombs.
\end{proof}
\begin{definition}[Collinear Three Charge Figure]
  \label{def:physics:phyx-mini-0900:phyxminiproblems-problemphyxmini0900-collinearthreechargefigure}
  \lean{PhyXMiniProblems.ProblemPhyXMini0900.CollinearThreeChargeFigure}
  \uses{def:physics:phyx-mini-0900:phyxminiproblems-problemphyxmini0900-chargesite, def:physics:phyx-mini-0900:phyxminiproblems-problemphyxmini0900-adjacentgap, def:physics:phyx-mini-0900:phyxminiproblems-problemphyxmini0900-figurechargecolor, def:physics:phyx-mini-0900:phyxminiproblems-problemphyxmini0900-figuresignglyph}
  Literal presentation data transcribed from primary image 900.png
\end{definition}
\begin{proof}
  This is the declared model definition of Collinear Three Charge Figure.
\end{proof}
\begin{definition}[Collinear Three Charge Equilibrium Setup]
  \label{def:physics:phyx-mini-0900:phyxminiproblems-problemphyxmini0900-collinearthreechargeequilibriumsetup}
  \lean{PhyXMiniProblems.ProblemPhyXMini0900.CollinearThreeChargeEquilibriumSetup}
  \uses{def:physics:phyx-mini-0900:phyxminiproblems-problemphyxmini0900-signedchargequantity, def:physics:phyx-mini-0900:phyxminiproblems-problemphyxmini0900-axialpositionquantity, def:physics:phyx-mini-0900:phyxminiproblems-problemphyxmini0900-lengthquantity, def:physics:phyx-mini-0900:phyxminiproblems-problemphyxmini0900-axialforcecomponentquantity, def:physics:phyx-mini-0900:phyxminiproblems-problemphyxmini0900-chargesite, def:physics:phyx-mini-0900:phyxminiproblems-problemphyxmini0900-adjacentgap, def:physics:phyx-mini-0900:phyxminiproblems-problemphyxmini0900-collinearthreechargefigure}
  Independent physical quantities for the three point charges. In particular, charge .q1 is not defined from an answer choice or from the requested value
\end{definition}
\begin{proof}
  This is the declared model definition of Collinear Three Charge Equilibrium Setup.
\end{proof}
\begin{definition}[displacement From Source To Q2 In Meters]
  \label{def:physics:phyx-mini-0900:phyxminiproblems-problemphyxmini0900-displacementfromsourcetoq2inmeters}
  \lean{PhyXMiniProblems.ProblemPhyXMini0900.displacementFromSourceToQ2InMeters}
  \uses{def:physics:phyx-mini-0900:phyxminiproblems-problemphyxmini0900-positioninmeters, def:physics:phyx-mini-0900:phyxminiproblems-problemphyxmini0900-chargesite, def:physics:phyx-mini-0900:phyxminiproblems-problemphyxmini0900-collinearthreechargeequilibriumsetup}
  Metre displacement from a source charge to the target charge q₂
\end{definition}
\begin{proof}
  This is the declared model definition of displacement From Source To Q2 In Meters.
\end{proof}
\begin{definition}[Matches Three Point Charge Equilibrium Scenario]
  \label{def:physics:phyx-mini-0900:phyxminiproblems-problemphyxmini0900-matchesthreepointchargeequilibriumscenario}
  \lean{PhyXMiniProblems.ProblemPhyXMini0900.MatchesThreePointChargeEquilibriumScenario}
  \uses{def:physics:phyx-mini-0900:phyxminiproblems-problemphyxmini0900-collinearthreechargeequilibriumsetup}
  The problem treats all three labelled circles as point charges and q₂ as the charge in static equilibrium
\end{definition}
\begin{proof}
  This is the declared model definition of Matches Three Point Charge Equilibrium Scenario.
\end{proof}
\begin{definition}[Matches Supplied Collinear Charge Figure]
  \label{def:physics:phyx-mini-0900:phyxminiproblems-problemphyxmini0900-matchessuppliedcollinearchargefigure}
  \lean{PhyXMiniProblems.ProblemPhyXMini0900.MatchesSuppliedCollinearChargeFigure}
  \uses{def:physics:phyx-mini-0900:phyxminiproblems-problemphyxmini0900-chargeinnanocoulombs, def:physics:phyx-mini-0900:phyxminiproblems-problemphyxmini0900-lengthinmeters, def:physics:phyx-mini-0900:phyxminiproblems-problemphyxmini0900-lengthincentimeters, def:physics:phyx-mini-0900:phyxminiproblems-problemphyxmini0900-positioninmeters, def:physics:phyx-mini-0900:phyxminiproblems-problemphyxmini0900-expectedchargetext, def:physics:phyx-mini-0900:phyxminiproblems-problemphyxmini0900-expectedchargecolor, def:physics:phyx-mini-0900:phyxminiproblems-problemphyxmini0900-expectedsignglyph, def:physics:phyx-mini-0900:phyxminiproblems-problemphyxmini0900-expectedprintedchargenanocoulombs, def:physics:phyx-mini-0900:phyxminiproblems-problemphyxmini0900-collinearthreechargeequilibriumsetup}
  All literal labels and geometric relations read from the primary raster. The coordinate origin is chosen at q₁; this convention does not constrain its charge. The only numerical charge premise is the middle +5.0 nC readout
\end{definition}
\begin{proof}
  This is the declared model definition of Matches Supplied Collinear Charge Figure.
\end{proof}
\begin{definition}[Has Physical Three Charge Parameters]
  \label{def:physics:phyx-mini-0900:phyxminiproblems-problemphyxmini0900-hasphysicalthreechargeparameters}
  \lean{PhyXMiniProblems.ProblemPhyXMini0900.HasPhysicalThreeChargeParameters}
  \uses{def:physics:phyx-mini-0900:phyxminiproblems-problemphyxmini0900-chargeincoulombs, def:physics:phyx-mini-0900:phyxminiproblems-problemphyxmini0900-lengthinmeters, def:physics:phyx-mini-0900:phyxminiproblems-problemphyxmini0900-collinearthreechargeequilibriumsetup}
  Nondegeneracy conditions implicit in a three-charge equilibrium problem
\end{definition}
\begin{proof}
  This is the declared model definition of Has Physical Three Charge Parameters.
\end{proof}
\begin{definition}[Satisfies Axial Point Charge Coulomb Force Law]
  \label{def:physics:phyx-mini-0900:phyxminiproblems-problemphyxmini0900-satisfiesaxialpointchargecoulombforcelaw}
  \lean{PhyXMiniProblems.ProblemPhyXMini0900.SatisfiesAxialPointChargeCoulombForceLaw}
  \uses{def:physics:phyx-mini-0900:phyxminiproblems-problemphyxmini0900-chargeincoulombs, def:physics:phyx-mini-0900:phyxminiproblems-problemphyxmini0900-forcecomponentinnewtons, def:physics:phyx-mini-0900:phyxminiproblems-problemphyxmini0900-collinearthreechargeequilibriumsetup, def:physics:phyx-mini-0900:phyxminiproblems-problemphyxmini0900-displacementfromsourcetoq2inmeters}
  For either other source, the signed force on q₂ is k q₂ q\_source (x₂ - x\_source) / |x₂ - x\_source|³. This is the one-dimensional vector form of Coulomb's law, stated only away from self-interaction. It contains no numerical conclusion for q₁
\end{definition}
\begin{proof}
  This is the declared model definition of Satisfies Axial Point Charge Coulomb Force Law.
\end{proof}
\begin{definition}[Satisfies Force Superposition On Q2]
  \label{def:physics:phyx-mini-0900:phyxminiproblems-problemphyxmini0900-satisfiesforcesuperpositiononq2}
  \lean{PhyXMiniProblems.ProblemPhyXMini0900.SatisfiesForceSuperpositionOnQ2}
  \uses{def:physics:phyx-mini-0900:phyxminiproblems-problemphyxmini0900-forcecomponentinnewtons, def:physics:phyx-mini-0900:phyxminiproblems-problemphyxmini0900-collinearthreechargeequilibriumsetup}
  The total horizontal force on q₂ is the sum of the two other charges' electrostatic forces
\end{definition}
\begin{proof}
  This is the declared model definition of Satisfies Force Superposition On Q2.
\end{proof}
\begin{definition}[Satisfies Static Equilibrium Force Balance]
  \label{def:physics:phyx-mini-0900:phyxminiproblems-problemphyxmini0900-satisfiesstaticequilibriumforcebalance}
  \lean{PhyXMiniProblems.ProblemPhyXMini0900.SatisfiesStaticEquilibriumForceBalance}
  \uses{def:physics:phyx-mini-0900:phyxminiproblems-problemphyxmini0900-forcecomponentinnewtons, def:physics:phyx-mini-0900:phyxminiproblems-problemphyxmini0900-collinearthreechargeequilibriumsetup}
  Static equilibrium entails a vanishing net horizontal force
\end{definition}
\begin{proof}
  This is the declared model definition of Satisfies Static Equilibrium Force Balance.
\end{proof}
\begin{lemma}[source Forces On Q2 sum to zero]
  \label{lem:physics:phyx-mini-0900:phyxminiproblems-problemphyxmini0900-sourceforcesonq2-sum-to-zero}
  \lean{PhyXMiniProblems.ProblemPhyXMini0900.sourceForcesOnQ2_sum_to_zero}
  \uses{def:physics:phyx-mini-0900:phyxminiproblems-problemphyxmini0900-forcecomponentinnewtons, def:physics:phyx-mini-0900:phyxminiproblems-problemphyxmini0900-collinearthreechargeequilibriumsetup, def:physics:phyx-mini-0900:phyxminiproblems-problemphyxmini0900-matchesthreepointchargeequilibriumscenario, def:physics:phyx-mini-0900:phyxminiproblems-problemphyxmini0900-satisfiesforcesuperpositiononq2, def:physics:phyx-mini-0900:phyxminiproblems-problemphyxmini0900-satisfiesstaticequilibriumforcebalance}
  The scenario and the two general force-combination laws first give a zero sum of the left-source and middle-source force components. This is a derived relation, not a premise about q₁
\end{lemma}
\begin{proof}
  The conclusion follows from the governing relations and figure readouts named in the statement of source Forces On Q2 sum to zero.
\end{proof}
\begin{definition}[Answer Choice]
  \label{def:physics:phyx-mini-0900:phyxminiproblems-problemphyxmini0900-answerchoice}
  \lean{PhyXMiniProblems.ProblemPhyXMini0900.AnswerChoice}
  Labels attached to the four displayed charge choices
\end{definition}
\begin{proof}
  This is the declared model definition of Answer Choice.
\end{proof}
\begin{definition}[charge In Nanocoulombs]
  \label{def:physics:phyx-mini-0900:phyxminiproblems-problemphyxmini0900-answerchoice-chargeinnanocoulombs}
  \lean{PhyXMiniProblems.ProblemPhyXMini0900.AnswerChoice.chargeInNanocoulombs}
  \uses{def:physics:phyx-mini-0900:phyxminiproblems-problemphyxmini0900-chargeinnanocoulombs, def:physics:phyx-mini-0900:phyxminiproblems-problemphyxmini0900-answerchoice}
  Nanocoulomb value printed next to each answer label
\end{definition}
\begin{proof}
  This is the declared model definition of charge In Nanocoulombs.
\end{proof}
\begin{definition}[recorded Dataset Answer]
  \label{def:physics:phyx-mini-0900:phyxminiproblems-problemphyxmini0900-recordeddatasetanswer}
  \lean{PhyXMiniProblems.ProblemPhyXMini0900.recordedDatasetAnswer}
  \uses{def:physics:phyx-mini-0900:phyxminiproblems-problemphyxmini0900-answerchoice}
  The answer label recorded in the supplied dataset metadata
\end{definition}
\begin{proof}
  This is the declared model definition of recorded Dataset Answer.
\end{proof}
% --- Archon declaration topology end ---
% --- Archon physics formalization source end ---
